\documentclass[12pt]{article}
\usepackage[utf8]{inputenc}
\usepackage[T1]{fontenc}
\usepackage[margin=1in]{geometry}
\usepackage{setspace}
\usepackage{xcolor}
\usepackage[colorlinks=true, urlcolor=blue!60!black, linkcolor=blue!60!black]{hyperref}
\setlength{\parindent}{2em}
\setlength{\parskip}{1.25em}
\renewcommand{\baselinestretch}{1.0}
\begin{document}
\begin{center}
{\LARGE \textit{Katzenbach v. Morgan}, 384 U.S. 641 (1966)\par}
\vspace{0.8em}
{\large Case Brief --- Four-Page Edition\par}
\end{center}
\vspace{1em}

\section{Memory Jogger}

The Supreme Court held that Section 5 of the Fourteenth Amendment is a positive grant of legislative power governed by \textit{McCulloch v. Maryland}, allowing Congress to override New York's English literacy test for Puerto Rican voters without a prior judicial finding that the test itself violated the Equal Protection Clause (\href{https://supreme.justia.com/cases/federal/us/384/641/}{Justia}).

\section{Detailed Case Facts}

New York's Constitution and Election Law required new voters to show they could read and write English, by passing a literacy test or proving completion of sixth grade in an English-language school. Because Puerto Rican schools taught predominantly in Spanish, the rule disenfranchised many of the several hundred thousand U.S. citizens who had moved from Puerto Rico to New York City. Senators Robert Kennedy and Jacob Javits added \S 4(e) to the Voting Rights Act of 1965 without committee hearings: no one who completed sixth grade in an accredited ``American-flag'' school in Puerto Rico could be denied the vote for inability to read English. In \textit{Lassiter v. Northampton County Board of Elections} (1959), the Court had held that a neutral literacy test does not in all circumstances violate the Constitution. New York City voters argued that, because New York's test had never been held unconstitutional, \S 4(e) exceeded Congress's \S 5 power and violated the Tenth Amendment.

\section{Procedural History}

The voters sued Attorney General Katzenbach in the District Court for the District of Columbia, which had exclusive jurisdiction under the Act. The New York City Board of Elections defended \S 4(e), the United States intervened, and New York's Attorney General attacked the law as amicus. A three-judge court held 2--1 (Judge McGowan dissenting) that Congress had exceeded its powers and enjoined enforcement. On direct appeal, argued with \textit{Cardona v. Power}, the Supreme Court reversed on June 13, 1966.

\section{Judicial Votes}

Reversed, \textbf{7--2}.
\begin{itemize}
\item \textbf{Majority:} Justice Brennan, joined by Chief Justice Warren and Justices Black, Clark, White, and Fortas, and by Justice Douglas in part.
\item \textbf{Partial concurrence:} Justice Douglas reserved the Fifth Amendment challenge to the American-flag limitation.
\item \textbf{Dissent:} Justice Harlan, joined by Justice Stewart.
\end{itemize}

\section{Holding}

Section 4(e) validly enforces the Fourteenth Amendment, and under the Supremacy Clause New York's literacy requirement yields to it. Section 5 power is not limited to striking state laws the courts have already held unconstitutional; legislation is ``appropriate'' under the \textit{McCulloch} standard if the Court can perceive a rational basis for Congress's judgment that it enforces equal protection. Brennan offered two alternative rationales: a \emph{remedial} one (the vote would help Puerto Ricans secure nondiscriminatory public services) and a \emph{substantive} one (Congress could itself judge the test invidious). Footnote 10's ``one-way ratchet''---Congress may expand but never dilute rights---was later disavowed in \textit{City of Boerne v. Flores} (1997), which kept only the remedial rationale and imposed ``congruence and proportionality.''

\section{Analysis of Opinions}

\textbf{Majority (Brennan).} Requiring a prior judicial finding would reduce Congress to ``the insignificant role of abrogating only those state laws that the judicial branch was prepared to adjudge unconstitutional.'' \textit{Lassiter} concerned judicial review under \S 1, not legislation under \S 5. Drawing on \textit{Ex parte Virginia} and \textit{South Carolina v. Katzenbach}, Brennan read \S 5 as broadly as the Necessary and Proper Clause. Congress could find that the ballot, the right ``preservative of all rights,'' would protect Puerto Ricans in schools, housing, and policing, or could weigh the test's nativist origins and Spanish-language media against New York's interests. Footnote 10 answered Harlan: Congress may not dilute guarantees. The American-flag limit was valid because reform need not ``strike at all evils at the same time.''

\textbf{Douglas (in part).} Joined the \S 5 holding but reserved the underinclusiveness question for an excluded litigant.

\textbf{Harlan, with Stewart (dissenting).} Under \textit{Lassiter}, New York's test was rational. Section 5 grants remedial power once a violation is shown, but ``it is a judicial question whether the condition with which Congress has thus sought to deal is, in truth, an infringement of the Constitution.'' Letting Congress define \S 1 would let it dilute rights too, and Congress built no record of discrimination, acting by ``mere \textit{ipse dixit}.''

\section{Examples: Future Applications}

\begin{enumerate}
\item \textbf{Same-Side: Courthouse Access.} After hearings documenting the exclusion of wheelchair users from courtrooms, Congress requires accessible courthouses and abrogates state immunity. \textit{Tennessee v. Lane} (2004) upheld this application of ADA Title II.
\item \textbf{Same-Side: Nationwide Literacy-Test Ban.} On a record that literacy tests perpetuate segregated education, Congress bans them nationwide; \textit{Oregon v. Mitchell} (1970) unanimously upheld the ban.
\item \textbf{Opposite-Side: Overriding a Free Exercise Ruling.} Congress imposes strict scrutiny on all state laws burdening religion without a record of bigotry. \textit{City of Boerne} struck down RFRA as applied to the States, adopting Harlan's dissent.
\item \textbf{Opposite-Side: Approving Segregation.} Congress declares equally funded segregated schools consistent with equal protection. Footnote 10 bars any power to ``restrict, abrogate, or dilute'' guarantees.
\item \textbf{Fence-Sitter: Voting Age of 18.} Congress lowers the voting age in state elections. \textit{Morgan}'s substantive theory supports it, but age gets only rational-basis review and States set qualifications; \textit{Oregon v. Mitchell} split 5--4, upholding it for federal elections only.
\end{enumerate}

\section{Critique}

\textbf{Originalist and textualist views.} Critics echo Harlan: the power to ``\textit{enforce}'' presupposes a fixed \S 1 baseline. Stephen Carter showed the framers rejected broader ``necessary and proper'' language and that the one-way ratchet fails when one group's asserted right collides with another's. On this view \textit{Morgan}'s substantive rationale bypasses Article V and \textit{Marbury}.

\textbf{Progressive and doctrinal views.} Robert Post and Reva Siegel defend \textit{Morgan} as democratic constitutionalism and call \textit{Boerne}'s limits ``juricentric,'' noting the 39th Congress meant \S 5 to empower Congress. Archibald Cox warned, however, that \textit{McCulloch} deference without a limiting principle leaves almost no judicial check. The case pits judicial supremacy and federalism against national civil rights enforcement.

\section{Key Quotations}

\begin{quote}
``Thus, the \textit{McCulloch v. Maryland} standard is the measure of what constitutes `appropriate legislation' under \S 5 of the Fourteenth Amendment. Correctly viewed, \S 5 is a positive grant of legislative power authorizing Congress to exercise its discretion in determining whether and what legislation is needed to secure the guarantees of the Fourteenth Amendment.''
\end{quote}
\noindent --- Justice Brennan, Opinion of the Court, 384 U.S. at 651.

\begin{quote}
``We emphasize that Congress' power under \S 5 is limited to adopting measures to enforce the guarantees of the Amendment; \S 5 grants Congress no power to restrict, abrogate, or dilute these guarantees.''
\end{quote}
\noindent --- Justice Brennan, Opinion of the Court, 384 U.S. at 651 n. 10.

\begin{quote}
``It is not for us to review the congressional resolution of these factors. It is enough that we be able to perceive a basis upon which the Congress might resolve the conflict as it did.''
\end{quote}
\noindent --- Justice Brennan, Opinion of the Court, 384 U.S. at 653.

\end{document}
